
\paragraph{M2: what the surrogate now does.} On the particle set alone, the set encoder beats the M1 trees on every
Tier 0 target (reconstructed: log loss 0.364 vs 0.423; MINOS match
0.060 vs 0.101; charge 0.119 vs 0.127)
and on the prong multiplicity (0.803 vs 0.843), so the raw
final state carries information that 48 hand-built summaries lose. The 9 generated Tier 1 variables close their
marginals to $W_1 \le 0.06$ in model units (0.005 for the matched-muon
momentum ratio, 0.003 for $\log E_\mathrm{recoil}$), the correlation matrix agrees to
$|\Delta\rho| \le 0.04$, and the muon response versus true momentum and the recoil versus hadronic
kinetic energy follow MasterAnaDev in median and width. The closure classifier on the reco vector alone reaches AUC
0.569 (blocks: muon 0.51, vertex 0.53, calorimetry 0.54); with the truth features added it
reaches 0.617, with the flags and multiplicity perfectly calibrated (0.49) and the residual
coming from the vertex block (0.58), then calorimetry (0.53) and muon (0.51). The M2 exit criterion of
$\mathrm{AUC} < 0.55$ is therefore met approximately for the marginal test and not for the conditional one.

\paragraph{What the residual is.} The reconstructed vertex $z$ is snapped to scintillator-plane positions for about
a third of the events: in a 10\,mm slice of true $z$ around 7000\,mm, 111 of 331 reconstructed vertices sit at
exactly 7007.0\,mm. Conditional on the truth the $z$ residual is therefore a comb, which a continuous flow smooths
out; the marginal is reproduced but the conditional structure is not, exactly what the two AUCs say. The natural
fix is a discrete head for the plane index with a continuous offset, the same pattern used for the multiplicity,
and is part of M3, the prong model, where discrete-plus-continuous outputs are needed anyway. Two further
limitations are documented: the two calibrated recoil variants are derived rather than generated
(Section~\ref{sec:m2-iter}), and the muon response in the 1--1.5\,GeV bin (3\% of events, dominated by unmatched
muons) has a median offset of about 0.7 model units between MasterAnaDev and the surrogate.

\paragraph{Lessons for the design.} Three findings generalise beyond M2. First, closure classifiers with block-wise
and conditional variants are far more informative than marginal plots: every defect found in this milestone was
invisible in the one-dimensional histograms. Second, the reco variables of an analysis tuple are not all
independent random variables; some are deterministic functions of others (calibrations), some are quantised
(plane positions, sentinels), and a surrogate has to represent that structure explicitly rather than hope a
density estimator discovers it. Third, sentinel conventions matter: $-9999.9$ for an undefined charge, $P=1,
E=0$ for a direction-only prong, NaN and $10^{14}$ for a handful of corrupt rows all had to be handled before the
statistics were meaningful.

\paragraph{Deliverable.} \texttt{scripts/surrogate\_to\_ntuple.py} runs the trained model on any truth-level
input with the interface of Section~\ref{sec:setup} and writes a pruned MasterAnaDev-style ROOT tree (46
branches: truth passthrough, Tier 0 flags, muon, vertex, calorimetry, multiplicity, and the per-event
reconstruction probability for reweighting). On 50\,000 truth events of the example file it reconstructs
39.6\%, as MasterAnaDev does. Prong-level branches follow with M3.
